\VerdierPhysicalPageAnchor{0081}%
\VerdierSemanticAnchor{ch1-2-5-11}
\noindent\textbf{2.5.11.}\quad
\VerdierSemanticAnchor{ch1-2-5-11-homotopism}
Let $X$ and $Y$ be two complexes of $C$. A morphism in
$\operatorname{comp}(C)$, $f\colon X\longrightarrow Y$, is called a
\emph{homotopy equivalence} if its image $\dot f$ in $K(C)$ is an
isomorphism.
\VerdierSemanticAnchor{ch1-2-5-11-homotopic-complexes}
The complexes $X$ and $Y$ are said to be \emph{homotopy equivalent} if they
are isomorphic in $K(C)$.

\VerdierSemanticAnchor{ch1-2-5-12}
\noindent\textbf{2.5.12.}\quad
\VerdierSemanticAnchor{ch1-2-5-12-setup}
Let $X$ and $Y$ be two complexes of $C$,
$f,g\colon X\rightrightarrows Y$ two morphisms in
$\operatorname{comp}(C)$, and
$a,b\colon f\rightsquigarrow g$ two homotopies connecting $f$ to $g$.
\VerdierSemanticAnchor{ch1-2-5-12-second-order-homotopy}
A \emph{second-order homotopy} connecting $a$ to $b$, denoted
$\alpha\colon a\rightsquigarrow b$, is a family of morphisms
\[
  \alpha=(\alpha^i\colon X^{i+2}\longrightarrow Y^i,\ i\in\mathbb Z)
\]
such that:

\VerdierSemanticAnchor{ch1-2-5-12-eq1}
\begin{equation}
  b^i-a^i
  =\alpha^i d_X^{i+1}-d_Y^{i-1}\alpha^{i-1}.
  \tag{2.5.12.1}
\end{equation}

\VerdierSemanticAnchor{ch1-2-5-13}
\noindent\textbf{2.5.13.}\quad
Let us note, to conclude, that the isomorphism (2.4.1)

\VerdierSemanticAnchor{ch1-2-5-13-complex-opposite-isomorphism}
\begin{equation*}
  \iota_1\colon
  \operatorname{comp}(C)^\circ
  \longrightarrow
  \operatorname{comp}(C^\circ)
\end{equation*}
defines, by virtue of (2.5.7), a unique isomorphism

\VerdierSemanticAnchor{ch1-2-5-13-homotopy-opposite-isomorphism}
\begin{equation*}
  \iota_1\colon
  K(C)^\circ
  \longrightarrow
  K(C^\circ)
\end{equation*}
\VerdierSemanticAnchor{ch1-2-5-13-z-category-isomorphism}
which is an isomorphism of $\mathbb Z$-categories.
\VerdierSemanticAnchor{ch1-2-5-13-identification}
Henceforth we shall most often identify the categories $K(C)^\circ$ and
$K(C^\circ)$ by this isomorphism.
